(1) 入射 Jones 向量为 $(1,1)^T/\sqrt2$，相位延迟片仅改变两正交分量的相对相位，故透过为 $\mathbf e=(1,e^{i\delta})^T/\sqrt2$。模平方 $|E_x|^2+|E_y|^2=1$，说明无吸收波片不会改变总强度。
(2) 按题中明确规定的 Stokes 符号约定，$S_0=1$、$S_1=0$、$S_2=\cos\delta$、$S_3=\sin\delta$。恒等式 $S_1^2+S_2^2+S_3^2=1=S_0^2$ 表明输出仍是完全偏振态；若换用另一种右左旋圆偏振约定，只有 $S_3$ 符号随之更换。
(3) 检偏器 Jones 投影为 $(\cos\beta,\sin\beta)\mathbf e=(\cos\beta+e^{i\delta}\sin\beta)/\sqrt2$。取模平方得 $I(\beta)/I_0=\frac12[1+\sin(2\beta)\cos\delta]$；改变偏振片角度只选择光强分量，不会使无损相位延迟片产生或消灭光子能量。
(4) $\delta=0$ 时仍是 $45^\circ$ 线偏振，$I/I_0=\cos^2(\beta-\pi/4)$；$\delta=\pi$ 时为 $-45^\circ$ 线偏振，$I/I_0=\cos^2(\beta+\pi/4)$。$\delta=\pi/2$ 时两分量等幅正交相位，为圆偏振，任意线检偏器均透过 $I_0/2$。三种情况下波片透过总强度均保持 $I_0$。
